\subsection{Erroneous behavior}

\begin{frame}[t,fragile]{Uninitialized variables}
	\begin{itemize}
		\item \textmark{Uninitialized variables}: Reading an uninitialized scalar (e.g. \texttt{int}) is undefined behavior.
		      \begin{lstlisting}
void f() {
  int count;
  ++count; // undefined behavior
  std::println("count = {}", count);
}
\end{lstlisting}

		      \mode<presentation>{\vfill\pause}
		\item \textmark{Traditional mitigations}:
		      \begin{itemize}
			      \item Static analysis tools (e.g., Clang-Tidy, PVS-Studio)
			      \item Compiler warnings (e.g., -Wuninitialized)
			      \item Code reviews and manual inspections
		      \end{itemize}
		      \begin{lstlisting}
void f() {
  int count = 0; // initialize count to a known value
  ++count; // defined behavior
  std::println("count = {}", count);
}
\end{lstlisting}
	\end{itemize}
\end{frame}

\begin{frame}[t,fragile]{C++26: Erroneous behavior}
	\begin{itemize}
		\item C++26 introduces the concept of \textmark{erroneous behavior}.
      \begin{itemize}
        \item Well-defined behavior that implementations are recommended to diagnose and report.
        \item A variable takes on a well-defined but implementation-defined value.
      \end{itemize}
\begin{lstlisting}
void f() {
  int count;
  ++count; // Erroneous behavior: count is initialized to an implementation-defined value
  std::println("count = {}", count);
}
\end{lstlisting}

    \mode<presentation>{\vfill\pause}
		\item Benefits of diagnosing erroneous behavior:
      \begin{itemize}
        \item Early detection of potential issues during development.
        \item Improved code quality and reliability.
        \item Encourages developers to write safer code by avoiding reliance on implementation-defined behavior.
      \end{itemize}
	\end{itemize}
\end{frame}

\begin{frame}[t,fragile]{Erroneous behavior: Implementation and diagnostics}
\begin{itemize}
  \item Implementations are encouraged to provide diagnostics for erroneous behavior.
    \begin{itemize}
      \item Compiler warnings or errors when detecting uninitialized variable usage.
      \item Runtime checks that can detect and report erroneous behavior during execution.
    \end{itemize}
\begin{lstlisting}
void f() {
  int count;
  ++count; // Compiler warning: 'count' is uninitialized
  std::println("count = {}", count);
}
\end{lstlisting}

    \mode<presentation>{\vfill\pause}
    \item An escape-hatch attribute for intentional non-initialization:
\begin{lstlisting}
void f() {
  [[indeterminate]] std::array<int, 1024> buffer; // Buffer is intentionally left uninitialized
  read_from_file(buffer); // Function that fills the buffer with data
  //...
}
\end{lstlisting}
\end{itemize}
\end{frame}